\documentclass{../../note}

\usepackage{xcolor}
\usepackage{longtable}
\usepackage{array}
\usepackage{listings}

\lstset{
  basicstyle=\ttfamily\small,
  keywordstyle=\color{blue},
  commentstyle=\color{green!50!black},
  stringstyle=\color{red!60!black},
  numbers=left,
  numberstyle=\tiny,
  numbersep=6pt,
  breaklines=true,
  frame=single
}

\title{HUFE 数据结构第二讲}
\author{isomoes}

\begin{document}

\maketitle

\section{本次课定位}

本次课对应 5 次课安排中的第 2 次课，时长 3 小时。本讲不再包含“数组和广义表”，而是把时间集中到\textbf{树和二叉树}，重点突破二叉树性质、遍历、线索二叉树与哈夫曼树。

这一讲比第一讲明显更难，因为题目不再只考“定义记忆”，而是开始要求同学把\textbf{结构性质、遍历规则和应用题计算}连起来理解。

\subsection{本次课目标}

\begin{itemize}
  \item 理解树、森林、结点的度、层次、高度、叶子结点等基本概念。
  \item 掌握二叉树的定义与常见特殊形态，尤其是满二叉树与完全二叉树。
  \item 熟练掌握二叉树的几个核心性质，并能用于选择题和计算题。
  \item 掌握二叉树的顺序存储、链式存储以及常见下标关系。
  \item 掌握先序、中序、后序、层序遍历的访问顺序及典型逆向判断题。
  \item 了解线索二叉树的基本思想，能区分“孩子指针”和“线索指针”。
  \item 掌握哈夫曼树构造过程和带权路径长度 WPL 的计算。
\end{itemize}

\subsection{考试导向提醒}

这一讲在试卷中属于\textbf{高频 + 中高难}内容，常见出题方式主要有四类：

\begin{itemize}
  \item \textbf{概念判断题}：如树与二叉树的区别、完全二叉树的定义、线索化后的指针含义。
  \item \textbf{性质计算题}：如由结点层数求最大结点数，由叶子数推度为 2 的结点数。
  \item \textbf{遍历分析题}：给定一棵树写遍历序列，或由两种遍历结果反推出树形。
  \item \textbf{应用构造题}：如哈夫曼树构造、WPL 计算、层序编号关系判断。
\end{itemize}

所以本讲复习不能只记“名词解释”，而要建立一条清晰主线：\textbf{概念 - 性质 - 存储 - 遍历 - 应用}。

\subsection{3小时建议节奏}

\begin{center}
\begin{tabular}{|c|c|p{9cm}|}
\hline
阶段 & 时间 & 主要内容 \\
\hline
第1段 & 35分钟 & 树的基本概念、术语、树与森林、树的常见表示思想 \\
\hline
第2段 & 55分钟 & 二叉树定义、特殊二叉树、核心性质、顺序与链式存储 \\
\hline
第3段 & 55分钟 & 先序、中序、后序、层序遍历，遍历序列分析题 \\
\hline
第4段 & 35分钟 & 线索二叉树、哈夫曼树、WPL 计算、课堂总结 \\
\hline
\end{tabular}
\end{center}

\section{树的基础概念}

\subsection{树的定义}

\textcolor{red}{树}是 $n$ 个结点的有限集合。当 $n=0$ 时，称为空树；当 $n>0$ 时，有且仅有一个特定结点称为\textbf{根结点}，其余结点可分成若干个互不相交的子集，每个子集本身又是一棵树，称为原树的\textbf{子树}。

这个定义里最重要的不是句子长，而是抓住三点：

\begin{itemize}
  \item 非空树有且仅有一个根。
  \item 除根外，其余结点都有且仅有一个直接前驱。
  \item 一个结点可以有多个直接后继，因此树是典型的一对多结构。
\end{itemize}

\subsection{树中的常见术语}

\begin{itemize}
  \item \textbf{结点的度}：该结点拥有的子树个数。
  \item \textbf{树的度}：整棵树中各结点度的最大值。
  \item \textbf{叶子结点}：度为 0 的结点，也叫终端结点。
  \item \textbf{分支结点}：度不为 0 的结点，也叫非终端结点。
  \item \textbf{孩子、双亲、兄弟}：由直接连接关系确定的相对关系。
  \item \textbf{路径}：从一个结点到另一个结点所经过的结点序列。
  \item \textbf{结点的层次}：从根开始定义，根通常记为第 1 层。
  \item \textbf{树的高度}：树中结点的最大层次，也叫树的深度。
\end{itemize}

\subsection{树与线性结构的区别}

第一讲中的线性表是一对一关系，而树是一对多关系，因此分析方法也会变化：

\begin{itemize}
  \item 线性结构更强调前驱和后继唯一。
  \item 树结构更强调层次、分支和递归定义。
  \item 在线性结构中常见的是“从前往后”；在树结构中常见的是“从根到子树”。
\end{itemize}

这也是为什么学树时，\textbf{递归思想}会变得很自然：因为“每棵子树仍然是树”。

\subsection{森林的概念}

\textbf{森林}是若干棵互不相交的树的集合。

考试中容易出现的理解是：

\begin{itemize}
  \item 一棵树去掉根结点后，可能变成一个森林。
  \item 给森林增加一个统一的根结点，可以转化为一棵树。
\end{itemize}

因此树和森林之间可以相互转化，只是观察角度不同。

\subsection{树的常见表示思想}

普通树不像二叉树那样天然只有两个孩子，所以表示方式更多是“了解性考点”：

\begin{enumerate}
  \item \textbf{双亲表示法}：每个结点记录其双亲位置，便于找父结点。
  \item \textbf{孩子表示法}：每个结点记录其所有孩子，便于找子结点。
  \item \textbf{孩子兄弟表示法}：每个结点用两个指针分别指向第一个孩子和下一个兄弟，可把普通树转成二叉链式形式。
\end{enumerate}

其中最值得记的是第三种，因为它为树和二叉树之间的联系提供了桥梁。

\subsection{树的基础练习}

\begin{enumerate}
  \item 树为什么属于非线性结构？
  \item 结点的度和树的度分别是什么意思？
  \item 叶子结点和分支结点如何区分？
  \item 一棵非空树中，除根结点外的其他结点各有什么共同特征？
  \item 请判断：森林中每个元素都必须是二叉树。（\hspace{1cm}）
\end{enumerate}

\section{二叉树}

\subsection{定义}

\textcolor{red}{二叉树}是另一类重要的树形结构。它的每个结点最多只有两棵子树，并且这两棵子树有左右之分，分别称为\textbf{左子树}和\textbf{右子树}。

这里要特别强调：二叉树不是简单的“度不超过 2 的树”，它还是\textbf{有序树}。

也就是说：

\begin{itemize}
  \item 左子树和右子树不能随意交换。
  \item 即使某结点只有一个孩子，也要区分它是左孩子还是右孩子。
\end{itemize}

\subsection{二叉树与普通树的区别}

\begin{enumerate}
  \item 二叉树每个结点最多两个孩子；普通树不一定。
  \item 二叉树左右子树有顺序；普通树的孩子次序未必是核心属性。
  \item 二叉树即使某层为空，也可能影响后续结构判断，尤其在完全二叉树中最明显。
\end{enumerate}

\subsection{特殊二叉树}

\subsubsection{满二叉树}

满二叉树指一棵高度为 $h$ 的二叉树，其每一层的结点数都达到最大值。此时总结点数为：

\[
2^h - 1
\]

它的特点是“长得最规整”，每个分支结点都有两个孩子，所有叶子都在同一层。

\subsubsection{完全二叉树}

完全二叉树指：除最后一层外，其余各层都满；最后一层的结点连续集中在最左边。

完全二叉树是考试重点，因为它和顺序存储联系最紧密。

常见判断角度：

\begin{itemize}
  \item 若最后一层有空位，只能出现在最右侧。
  \item 一旦某位置空了，它右边不能再出现结点。
  \item 完全二叉树不一定是满二叉树，但满二叉树一定是完全二叉树。
\end{itemize}

\subsubsection{斜树}

如果一棵二叉树所有结点都只有左孩子，称为左斜树；都只有右孩子，称为右斜树。

这种结构在遍历题中很容易判断，但在性质题中要注意：它虽然仍是二叉树，却已经很接近线性结构。

\subsection{二叉树的核心性质}

\subsubsection{性质1：第 $i$ 层最多有多少个结点}

二叉树第 $i$ 层上至多有：

\[
2^{i-1} \quad (i \geq 1)
\]

理解方法很简单：根第 1 层最多 1 个，第 2 层最多 2 个，第 3 层最多 4 个，每下一层都最多翻倍。

\subsubsection{性质2：高度为 $h$ 的二叉树最多有多少个结点}

若二叉树高度为 $h$，则最多有：

\[
2^h - 1
\]

这是把每一层最多结点数相加得到的结果：

\[
1 + 2 + 4 + \cdots + 2^{h-1} = 2^h - 1
\]

\subsubsection{性质3：度为 0 和度为 2 的结点关系}

对任意一棵非空二叉树，若叶子结点个数为 $n_0$，度为 2 的结点个数为 $n_2$，则有：

\[
n_0 = n_2 + 1
\]

这个性质几乎是必考性质之一。注意它只和\textbf{叶子结点数、度为 2 的结点数}直接相关，与度为 1 的结点数无直接等式关系。

\subsubsection{性质4：完全二叉树编号关系}

若完全二叉树按层序从 1 开始编号，则编号为 $i$ 的结点：

\begin{itemize}
  \item 双亲编号为 $\lfloor i/2 \rfloor$（当 $i > 1$）
  \item 左孩子编号为 $2i$
  \item 右孩子编号为 $2i+1$
\end{itemize}

这个关系是顺序存储判断题的核心基础。

\subsection{性质理解例子}

\begin{enumerate}
  \item 若某二叉树第 4 层，则该层最多有 $2^3 = 8$ 个结点。
  \item 若某二叉树高度为 5，则最多有 $2^5 - 1 = 31$ 个结点。
  \item 若某二叉树有 10 个叶子结点，则度为 2 的结点数一定为 9。
  \item 在完全二叉树中，若某结点编号为 6，则其双亲编号为 3，左孩子编号为 12，右孩子编号为 13。
\end{enumerate}

\subsection{二叉树的存储结构}

\subsubsection{顺序存储}

顺序存储通常适用于\textbf{完全二叉树}，因为它层次紧凑，不会浪费太多空间。

其核心思想是：按层序把结点依次放入数组。

优点：
\begin{itemize}
  \item 结点之间的位置关系可直接通过下标计算。
  \item 访问编号结点方便。
\end{itemize}

缺点：
\begin{itemize}
  \item 若不是完全二叉树，会产生较多空位，空间利用率下降。
\end{itemize}

\subsubsection{链式存储}

链式存储更适合一般二叉树，通常使用二叉链表：

\begin{lstlisting}[language=C]
typedef struct BiNode {
    ElemType data;
    struct BiNode *lchild;
    struct BiNode *rchild;
} BiNode, *BiTree;
\end{lstlisting}

这里：

\begin{itemize}
  \item \texttt{data} 存放结点值；
  \item \texttt{lchild} 指向左孩子；
  \item \texttt{rchild} 指向右孩子。
\end{itemize}

若某个孩子不存在，对应指针域为空。这些“空指针域”也是后面线索二叉树能够出现的基础。

\subsection{二叉树部分常见判断}

\begin{enumerate}
  \item 二叉树并不等价于度不超过 2 的有序树，这个说法虽然接近，但答题时仍应强调“左右子树有次序”。
  \item 满二叉树一定是完全二叉树；完全二叉树不一定是满二叉树。
  \item 只有完全二叉树的顺序存储下标关系才最自然、最常考。
  \item 二叉树中某结点只有一个孩子时，要区分它到底是左孩子还是右孩子。
\end{enumerate}

\subsection{二叉树练习}

\begin{enumerate}
  \item 二叉树为什么说是有序树？
  \item 满二叉树和完全二叉树有什么关系？
  \item 二叉树第 $i$ 层最多有多少个结点？
  \item 若某非空二叉树叶子结点数为 7，则度为 2 的结点数是多少？
  \item 在完全二叉树中，若结点编号为 9，则其双亲编号是多少？
  \item 请判断：完全二叉树的最后一层结点可以任意分散出现。（\hspace{1cm}）
\end{enumerate}

\section{二叉树的遍历}

\subsection{为什么遍历是重点}

遍历就是按某种规则，把树中所有结点不重不漏地访问一次。

考试中之所以高频考遍历，是因为它同时考了三件事：

\begin{itemize}
  \item 你是否真正理解“根、左、右”的递归结构；
  \item 你是否能区分不同遍历顺序；
  \item 你是否能由遍历序列反推树形或分析应用场景。
\end{itemize}

\subsection{先序遍历}

先序遍历的规则是：

\[
\text{根} \rightarrow \text{左子树} \rightarrow \text{右子树}
\]

记忆口诀：\textbf{先看根}。

\begin{lstlisting}[language=C]
void PreOrder(BiTree T) {
    if (T != NULL) {
        visit(T);
        PreOrder(T->lchild);
        PreOrder(T->rchild);
    }
}
\end{lstlisting}

先序遍历常用于“先处理当前结点，再处理子问题”的场景。

\subsection{中序遍历}

中序遍历的规则是：

\[
\text{左子树} \rightarrow \text{根} \rightarrow \text{右子树}
\]

记忆口诀：\textbf{根在中间}。

\begin{lstlisting}[language=C]
void InOrder(BiTree T) {
    if (T != NULL) {
        InOrder(T->lchild);
        visit(T);
        InOrder(T->rchild);
    }
}
\end{lstlisting}

中序遍历在二叉排序树中尤其重要，因为它会得到一个递增序列。

\subsection{后序遍历}

后序遍历的规则是：

\[
\text{左子树} \rightarrow \text{右子树} \rightarrow \text{根}
\]

记忆口诀：\textbf{最后看根}。

\begin{lstlisting}[language=C]
void PostOrder(BiTree T) {
    if (T != NULL) {
        PostOrder(T->lchild);
        PostOrder(T->rchild);
        visit(T);
    }
}
\end{lstlisting}

后序遍历常用于“先处理左右子问题，再处理当前结点”的场景，例如求树高、释放结点空间等。

\subsection{层序遍历}

层序遍历的规则是：按照从上到下、从左到右的顺序逐层访问。

它和前三种递归遍历不同，通常借助\textbf{队列}实现，因此也和第一讲内容形成了连接。

\begin{lstlisting}[language=C]
void LevelOrder(BiTree T) {
    InitQueue(Q);
    if (T != NULL)
        EnQueue(Q, T);
    while (!QueueEmpty(Q)) {
        DeQueue(Q, p);
        visit(p);
        if (p->lchild != NULL)
            EnQueue(Q, p->lchild);
        if (p->rchild != NULL)
            EnQueue(Q, p->rchild);
    }
}
\end{lstlisting}

\subsection{四种遍历的对比}

\begin{center}
\begin{tabular}{|c|c|c|}
\hline
遍历方式 & 访问顺序核心 & 记忆点 \\
\hline
先序 & 根左右 & 先看根 \\
\hline
中序 & 左根右 & 根在中间 \\
\hline
后序 & 左右根 & 最后看根 \\
\hline
层序 & 自上而下逐层 & 借助队列 \\
\hline
\end{tabular}
\end{center}

\subsection{遍历序列分析例子}

设某二叉树结构为：根结点为 $A$，左子树根为 $B$，右子树根为 $C$；其中 $B$ 的左右孩子分别为 $D$、$E$，$C$ 只有右孩子 $F$。

则有：

\begin{itemize}
  \item 先序遍历：$A, B, D, E, C, F$
  \item 中序遍历：$D, B, E, A, C, F$
  \item 后序遍历：$D, E, B, F, C, A$
  \item 层序遍历：$A, B, C, D, E, F$
\end{itemize}

这组序列非常适合反复练习，因为它包含了“左子树完整、右子树不完整”的典型情况。

\subsection{由遍历反推树形时要注意什么}

\begin{enumerate}
  \item 仅知道一种遍历序列，通常不能唯一确定一棵二叉树。
  \item 若已知\textbf{先序 + 中序}，通常可以唯一确定一棵二叉树。
  \item 若已知\textbf{后序 + 中序}，通常也可以唯一确定一棵二叉树。
  \item 若只有先序和后序，在一般情况下不能唯一确定。
\end{enumerate}

这里的核心原因是：中序遍历把“根左边”和“根右边”分得很清楚，因此能提供结构划分依据。

\subsection{遍历练习}

\begin{enumerate}
  \item 先序、中序、后序遍历分别把根结点放在什么位置访问？
  \item 层序遍历为什么通常要借助队列？
  \item 为什么中序遍历在反推二叉树时特别关键？
  \item 若某二叉树先序为 $A,B,D,E,C$，中序为 $D,B,E,A,C$，则根结点是谁？
  \item 请判断：只知道一棵二叉树的先序遍历序列，就一定能唯一确定该树。（\hspace{1cm}）
\end{enumerate}

\section{线索二叉树}

\subsection{为什么需要线索二叉树}

普通二叉链表中，很多结点的左指针或右指针可能为空。若这些空指针不利用，就会造成一定浪费。

线索二叉树的基本思想就是：把这些空指针改造为指向某种遍历次序下的\textbf{前驱}或\textbf{后继}。

因此它的目标不是“改变树”，而是让遍历和查找前驱后继更方便。

\subsection{基本概念}

若某结点原来的左指针为空，可让它指向在某种遍历序列中的前驱；若右指针为空，可让它指向该结点在该遍历序列中的后继。这样的指针就称为\textbf{线索}。

按照建立线索所依据的遍历方式不同，可分为：

\begin{itemize}
  \item 前序线索二叉树
  \item 中序线索二叉树
  \item 后序线索二叉树
\end{itemize}

其中考试中最常见的是\textbf{中序线索二叉树}。

\subsection{存储结构中的标志位}

为了区分一个指针到底是“孩子指针”还是“线索指针”，通常增加两个标志位：

\begin{lstlisting}[language=C]
typedef struct ThreadNode {
    ElemType data;
    struct ThreadNode *lchild, *rchild;
    int ltag, rtag;
} ThreadNode, *ThreadTree;
\end{lstlisting}

常见约定：

\begin{itemize}
  \item \texttt{ltag = 0} 表示 \texttt{lchild} 指向左孩子；\texttt{ltag = 1} 表示指向前驱线索。
  \item \texttt{rtag = 0} 表示 \texttt{rchild} 指向右孩子；\texttt{rtag = 1} 表示指向后继线索。
\end{itemize}

这类题目的关键不是背代码，而是分清“空指针被改造成线索以后，必须有标志位说明其意义”。

\subsection{线索二叉树易错点}

\begin{enumerate}
  \item 线索不是额外新加一条边，而是对原空指针域的利用。
  \item 建线索所依据的遍历不同，前驱和后继关系也会不同。
  \item 线索化后，不是所有指针都变成线索，原本存在孩子的指针仍然是孩子指针。
\end{enumerate}

\subsection{线索二叉树练习}

\begin{enumerate}
  \item 线索二叉树为什么能够提高遍历效率？
  \item 为什么线索二叉树需要增加标志位？
  \item 中序线索二叉树中的线索关系是依据哪一种遍历建立的？
  \item 请判断：线索二叉树中所有左右指针都一定改成线索。（\hspace{1cm}）
\end{enumerate}

\section{哈夫曼树}

\subsection{哈夫曼树的核心思想}

\textcolor{red}{哈夫曼树}也叫\textbf{最优二叉树}。它研究的不是“形状是否好看”，而是：在带权叶子结点已知时，如何使带权路径长度最小。

它最常考的两个点是：

\begin{itemize}
  \item 哈夫曼树如何构造；
  \item 带权路径长度 WPL 如何计算。
\end{itemize}

\subsection{路径长度与带权路径长度}

\begin{itemize}
  \item \textbf{路径长度}：从根到该结点所经过的边数。
  \item \textbf{带权路径长度}：结点权值乘以该结点路径长度。
  \item \textbf{WPL}：所有叶子结点带权路径长度之和。
\end{itemize}

若叶子结点权值分别为 $w_1, w_2, \dots, w_n$，对应路径长度为 $l_1, l_2, \dots, l_n$，则：

\[
WPL = \sum_{i=1}^{n} w_i l_i
\]

\subsection{哈夫曼树构造规则}

构造过程可概括为：

\begin{enumerate}
  \item 把每个权值看作一棵只有一个结点的树。
  \item 每次选出权值最小的两棵树。
  \item 把它们合并成一棵新树，新根权值等于两者之和。
  \item 重复上述过程，直到森林中只剩一棵树。
\end{enumerate}

最关键的规则就一句话：\textbf{每次取最小的两个合并}。

\subsection{哈夫曼树计算示例}

设权值集合为：

\[
\{2, 4, 5, 7\}
\]

构造步骤：

\begin{enumerate}
  \item 先取最小的 $2$ 和 $4$，合并成 $6$。
  \item 当前集合变为 $\{5, 6, 7\}$，再取 $5$ 和 $6$，合并成 $11$。
  \item 当前集合变为 $\{7, 11\}$，再合并成 $18$。
\end{enumerate}

若最终树中各叶子的路径长度分别为：

\begin{itemize}
  \item 权值 7 的路径长度为 1
  \item 权值 5 的路径长度为 2
  \item 权值 2 的路径长度为 3
  \item 权值 4 的路径长度为 3
\end{itemize}

则：

\[
WPL = 7 \times 1 + 5 \times 2 + 2 \times 3 + 4 \times 3 = 35
\]

\subsection{哈夫曼树的几个结论}

\begin{enumerate}
  \item 权值越大的叶子，通常离根越近，编码也越短。
  \item 哈夫曼树中没有度为 1 的结点。
  \item 哈夫曼树不一定唯一，但最小 WPL 相同。
\end{enumerate}

\subsection{哈夫曼编码的理解}

哈夫曼编码就是在哈夫曼树基础上形成的前缀编码。通常约定：

\begin{itemize}
  \item 向左走记为 0
  \item 向右走记为 1
\end{itemize}

这样从根走到某叶子形成的一串 0/1，就是该字符的编码。

它的优势在于：高频字符编码短，低频字符编码长，总体压缩效率更高。

\subsection{哈夫曼树练习}

\begin{enumerate}
  \item 哈夫曼树为什么又叫最优二叉树？
  \item WPL 的计算对象是所有结点还是所有叶子结点？
  \item 构造哈夫曼树时，每一步为什么都要选权值最小的两棵树？
  \item 请判断：哈夫曼树中可能存在只有一个孩子的结点。（\hspace{1cm}）
\end{enumerate}

\section{易混点总结}

\begin{enumerate}
  \item \textbf{树与二叉树}：二叉树不是普通树的特例描述那么简单，它强调左右次序。
  \item \textbf{满二叉树与完全二叉树}：满二叉树更严格，完全二叉树更常考。
  \item \textbf{第 $i$ 层最多结点数与高度为 $h$ 的最多结点数}：一个看单层，一个看整树。
  \item \textbf{先序、中序、后序}：关键不是背名字，而是根结点出现在访问序列中的位置。
  \item \textbf{层序遍历与前面三种遍历}：层序通常借助队列，不完全依赖递归。
  \item \textbf{线索与孩子}：线索二叉树中的指针不一定都指向孩子，需要靠标志位区分。
  \item \textbf{哈夫曼树与普通二叉树}：哈夫曼树关注的是 WPL 最小，而不是结点值大小规律。
\end{enumerate}

\section{本讲知识框架}

建议把本讲内容整理成下面这条主线：

\begin{center}
\begin{tabular}{|c|c|c|c|}
\hline
主题 & 核心问题 & 典型内容 & 高频考点 \\
\hline
树基础 & 什么是层次结构 & 根、孩子、度、森林 & 术语辨析 \\
\hline
二叉树 & 为什么左右有序 & 满二叉树、完全二叉树 & 性质计算、编号关系 \\
\hline
遍历 & 如何不重不漏访问 & 先中后层序 & 序列判断、反推结构 \\
\hline
线索化 & 空指针如何利用 & 前驱、后继、标志位 & 线索含义辨析 \\
\hline
哈夫曼树 & 怎样让 WPL 最小 & 构造规则、编码 & 合并步骤、WPL \\
\hline
\end{tabular}
\end{center}

\section{典型例题}

\subsection{例题1：二叉树第 $i$ 层结点数}

\textbf{题目：}某二叉树第 5 层最多有多少个结点？

\textbf{分析：}利用性质“第 $i$ 层最多有 $2^{i-1}$ 个结点”。

\textbf{答案：}$2^4 = 16$。

\subsection{例题2：叶子结点与度为 2 结点关系}

\textbf{题目：}某非空二叉树中叶子结点数为 12，则度为 2 的结点数是多少？

\textbf{分析：}由性质 $n_0 = n_2 + 1$，得 $n_2 = 11$。

\textbf{答案：}11。

\subsection{例题3：完全二叉树编号}

\textbf{题目：}在一棵按层序从 1 开始编号的完全二叉树中，编号为 10 的结点的双亲编号是多少？

\textbf{分析：}双亲编号为 $\lfloor 10/2 \rfloor$。

\textbf{答案：}5。

\subsection{例题4：遍历序列判断}

\textbf{题目：}某二叉树先序遍历为 $A,B,D,E,C,F$，中序遍历为 $D,B,E,A,C,F$，则根结点是谁？

\textbf{分析：}先序遍历第一个访问的一定是根结点。

\textbf{答案：}$A$。

\subsection{例题5：层序遍历使用什么结构}

\textbf{题目：}二叉树层序遍历通常借助什么线性结构实现？

\textbf{分析：}层序遍历要求先访问先进入下一层的结点，符合先进先出。

\textbf{答案：}队列。

\subsection{例题6：哈夫曼树的 WPL}

\textbf{题目：}一棵哈夫曼树的叶子权值和路径长度分别为 $(3,2)$、$(5,2)$、$(7,1)$，求其 WPL。

\textbf{解：}
\[
WPL = 3 \times 2 + 5 \times 2 + 7 \times 1 = 23
\]

\textbf{答案：}23。

\section{课堂练习}

\subsection{基础练习}

\begin{enumerate}
  \item 树为什么可以用递归思想来理解？
  \item 二叉树和普通树相比，最关键的差异是什么？
  \item 满二叉树与完全二叉树分别如何定义？
  \item 顺序存储为什么更适合完全二叉树？
  \item 先序、中序、后序遍历的根结点分别在什么时机访问？
\end{enumerate}

\subsection{判断与选择}

\begin{enumerate}
  \item 二叉树中某结点的左右子树可以随意交换而不影响结构定义。（\hspace{1cm}）
  \item 满二叉树一定是完全二叉树。（\hspace{1cm}）
  \item 完全二叉树适合采用顺序存储结构。（\hspace{1cm}）
  \item 只知道一棵二叉树的中序遍历序列，就一定能唯一确定该树。（\hspace{1cm}）
  \item 线索二叉树中的空指针域可能被用来存放前驱或后继信息。（\hspace{1cm}）
  \item 哈夫曼树的目标是让所有叶子尽量处在同一层。（\hspace{1cm}）
\end{enumerate}

\subsection{计算与分析}

\begin{enumerate}
  \item 某二叉树第 6 层最多有多少个结点？
  \item 若某二叉树高度为 4，则最多有多少个结点？
  \item 某非空二叉树有 15 个叶子结点，则度为 2 的结点数是多少？
  \item 在完全二叉树中，若某结点编号为 14，请写出其双亲、左孩子、右孩子的编号。
  \item 设某二叉树先序遍历为 $A,B,C,D$，后序遍历为 $D,C,B,A$，为什么一般不能只靠这两种序列唯一确定一棵二叉树？
  \item 若哈夫曼树叶子权值分别为 1、3、6，且其路径长度分别为 2、2、1，求 WPL。
\end{enumerate}

\subsection{综合小题}

\textbf{题目：}某系统需要完成两个任务：

\begin{itemize}
  \item 任务 A：按层次从上到下输出一棵二叉树的结点；
  \item 任务 B：根据若干字符出现频率生成压缩编码方案。
\end{itemize}

请分别指出任务 A 和任务 B 对应的核心数据结构思想，并说明理由。

\section{课后小测参考答案}

\begin{enumerate}
  \item 因为树中每棵子树本身仍然是树，整体与局部结构相同。
  \item 二叉树强调每个结点最多两个孩子，且左右子树有顺序。
  \item 满二叉树每层都达到最大结点数；完全二叉树除最后一层外都满，最后一层结点连续集中在最左边。
  \item 因为完全二叉树结构紧凑，按层序编号后可直接通过数组下标反映双亲和孩子关系。
  \item 先序先访问根，中序中间访问根，后序最后访问根。
  \item 第1题错；第2题对；第3题对；第4题错；第5题对；第6题错。
  \item 第 6 层最多有 $2^5 = 32$ 个结点。
  \item 高度为 4 时，最多有 $2^4 - 1 = 15$ 个结点。
  \item 由 $n_0 = n_2 + 1$ 得度为 2 的结点数为 14。
  \item 双亲为 7，左孩子为 28，右孩子为 29。
  \item 因为没有中序序列提供左右子树的划分边界，一般无法唯一确定结构。
  \item $WPL = 1 \times 2 + 3 \times 2 + 6 \times 1 = 14$。
\end{enumerate}

综合小题参考：任务 A 对应层序遍历，核心借助队列；任务 B 对应哈夫曼树，核心是按频率构造最优二叉树以得到较短编码。

\subsection{分节练习参考答案}

\textbf{树的基础练习参考：}

\begin{enumerate}
  \item 因为一个结点可以对应多个后继，元素之间不是一对一关系。
  \item 结点的度是某结点孩子个数；树的度是整棵树中结点度的最大值。
  \item 度为 0 的是叶子结点，度不为 0 的是分支结点。
  \item 都有且仅有一个直接前驱，也就是双亲结点。
  \item 错。森林中的元素是树，不要求每棵树都是二叉树。
\end{enumerate}

\textbf{二叉树练习参考：}

\begin{enumerate}
  \item 因为左右子树有明确区分，不能随意交换。
  \item 满二叉树一定是完全二叉树，但完全二叉树不一定是满二叉树。
  \item 第 $i$ 层最多有 $2^{i-1}$ 个结点。
  \item 6。
  \item 双亲编号为 $\lfloor 9/2 \rfloor = 4$。
  \item 错。最后一层必须连续集中在最左边。
\end{enumerate}

\textbf{遍历练习参考：}

\begin{enumerate}
  \item 先序先访问根；中序根在中间；后序最后访问根。
  \item 因为层序遍历要求先访问先进入队列的上一层结点，符合先进先出。
  \item 因为中序序列可以把根左边和根右边的结点范围划分出来。
  \item 根结点是 $A$。
  \item 错。只知道先序遍历通常不能唯一确定一棵二叉树。
\end{enumerate}

\textbf{线索二叉树练习参考：}

\begin{enumerate}
  \item 因为空指针可被利用来直接指向某种遍历次序下的前驱或后继。
  \item 因为同一个指针域可能表示孩子，也可能表示线索，必须区分。
  \item 依据中序遍历建立。
  \item 错。原本存在孩子的指针仍然指向孩子。
\end{enumerate}

\textbf{哈夫曼树练习参考：}

\begin{enumerate}
  \item 因为它使带权路径长度 WPL 达到最小。
  \item 是所有叶子结点。
  \item 因为这样才能保证最终带权路径长度最小。
  \item 错。哈夫曼树中不存在度为 1 的结点。
\end{enumerate}

\section{本讲小结}

本讲的核心不只是“认识树”，而是建立树形结构的分析方法：\textbf{先分清基本概念，再掌握二叉树性质，然后通过遍历理解递归结构，最后把线索二叉树和哈夫曼树看作具体应用}。

如果说第一讲是在建立数据结构的入门视角，那么这一讲就是正式进入“结构分析题”。其中最需要反复练习的三个点分别是：\textbf{完全二叉树编号关系、四种遍历顺序、哈夫曼树的构造与 WPL 计算}。

下一讲将进入图、查找和排序，计算量和综合性会进一步提高，因此这一讲的遍历与结构理解一定要打牢。

\end{document}
